%=============================================================================
% APPENDIX J: DERIVATION OF THE ψ-CMB SOLUTION
% Complete derivation of peak ratio and peak location from ψ-physics
%=============================================================================

\section{Derivation of the $\psi$-CMB Solution}\label{app:psi-cmb}

This appendix provides complete derivations of the $\psi$-CMB results presented in \S\ref{sec:psi-cmb}. We derive both the peak ratio $R \approx 2.34$ from baryon loading in $\psi$-gravity and the peak location $\ell_1 \approx 220$ from $\psi$-lensing.

%-----------------------------------------------------------------------------
\subsection{The $\psi$-Acoustic Oscillator}\label{app:psi-acoustic}
%-----------------------------------------------------------------------------

\paragraph{Setup.}
Consider a baryon-photon fluid in $\psi$-gravity. The temperature perturbation $\Theta \equiv \delta T/T$ obeys:
\begin{equation}\label{eq:app-acoustic}
\ddot{\Theta} + c_s^2(\psi) k^2 \Theta = -\frac{k^2}{1+R_b} \Phi_\psi,
\end{equation}
where:
\begin{itemize}[nosep]
\item $c_s(\psi) = c(\psi)/\sqrt{3}$ is the sound speed with $c(\psi) = c_0 e^{-\psi}$
\item $R_b = 3\rho_b/(4\rho_\gamma) \approx 0.6$ is the baryon-to-photon density ratio
\item $\Phi_\psi = \Phi/\mu(x)$ is the $\psi$-enhanced gravitational potential
\end{itemize}

\paragraph{Solution structure.}
The general solution has the form:
\begin{equation}
\Theta(k,\tau) = A(k) \cos(k r_s) + B(k) \sin(k r_s) + \text{(driving term)},
\end{equation}
where $r_s = \int c_s(\psi)\,d\tau$ is the sound horizon.

\paragraph{Peak/trough pattern.}
\begin{itemize}[nosep]
\item Odd peaks ($n = 1, 3, 5, \ldots$): compressions (maxima of $|\Theta|$)
\item Even peaks ($n = 2, 4, 6, \ldots$): rarefactions (minima of $|\Theta|$)
\end{itemize}

In standard cosmology, baryon loading causes compressions to be enhanced relative to rarefactions, producing the odd/even asymmetry.

%-----------------------------------------------------------------------------
\subsection{Peak Height Asymmetry}\label{app:peak-asymmetry}
%-----------------------------------------------------------------------------

\paragraph{The asymmetry factor.}
The ratio of odd to even peak heights is determined by the asymmetry factor $A$:
\begin{equation}
\frac{H_{\rm odd}}{H_{\rm even}} = \left(\frac{1+A}{1-A}\right).
\end{equation}

\paragraph{Factor decomposition.}
We decompose $A$ into four physically distinct contributions:
\begin{equation}\label{eq:app-A-decomp}
A = f_{\rm baryon} \times f_{\rm ISW} \times f_{\rm vis} \times f_{\rm Dop}.
\end{equation}

\subsubsection{Baryon Loading Factor $f_{\rm baryon}$}

The baryon-photon oscillator with baryon loading $R_b$ produces asymmetry:
\begin{equation}
f_{\rm baryon} = \frac{R_b}{\sqrt{1+R_b}}.
\end{equation}

\paragraph{Derivation.}
In the tight-coupling limit, the photon-baryon fluid satisfies:
\begin{equation}
\ddot{\Theta} + \frac{R_b}{1+R_b} \dot{\Theta} + \frac{c_s^2 k^2}{(1+R_b)} \Theta = -\frac{k^2 \Phi}{(1+R_b)}.
\end{equation}
The baryon drag term $\frac{R_b}{1+R_b}\dot{\Theta}$ introduces phase shift and amplitude modulation. For adiabatic perturbations with $\Phi = \text{const}$, the equilibrium compression is:
\begin{equation}
\Theta_{\rm eq} = -\Phi/(1+R_b).
\end{equation}
Oscillations about this equilibrium have amplitude modulated by $1/\sqrt{1+R_b}$. The asymmetry between compression (toward $\Theta_{\rm eq}$) and rarefaction (away from $\Theta_{\rm eq}$) gives:
\begin{equation}
f_{\rm baryon} = \frac{|\Theta_{\rm eq}|}{1/\sqrt{1+R_b}} = \frac{R_b}{\sqrt{1+R_b}}.
\end{equation}

\paragraph{Numerical value.}
With $R_b = 0.6$ (from BBN):
\begin{equation}
f_{\rm baryon} = \frac{0.6}{\sqrt{1.6}} = \frac{0.6}{1.265} = 0.474.
\end{equation}

\subsubsection{Integrated Sachs-Wolfe Factor $f_{\rm ISW}$}

The observed temperature perturbation includes the Sachs-Wolfe and integrated Sachs-Wolfe terms:
\begin{equation}
\frac{\Delta T}{T} = \Theta + \Phi + 2\int \dot{\Phi}\,d\tau.
\end{equation}

\paragraph{$\psi$-ISW effect.}
In $\psi$-gravity, the potential $\Phi_\psi = \Phi/\mu$ evolves as $\mu$ changes. If $\mu$ increases with time (gravity ``turns on''), $\Phi_\psi$ decays, producing an ISW contribution.

\paragraph{Cancellation.}
The SW term ($\Phi$) and ISW term ($2\int\dot{\Phi}\,d\tau$) partially cancel. In $\psi$-cosmology, this cancellation is approximately 50\%:
\begin{equation}
f_{\rm ISW} \approx 0.50.
\end{equation}

This value depends on the detailed $\mu$-evolution but is constrained to be $\mathcal{O}(0.5)$ by physical considerations.

\subsubsection{Visibility Function Factor $f_{\rm vis}$}

Recombination is not instantaneous. The visibility function $g(\tau) = \dot{\tau}_c e^{-\tau_c}$ has finite width $\Delta\tau$.

\paragraph{Effect on asymmetry.}
Finite-width recombination smears out the sharp features in the angular power spectrum. The effect on the asymmetry is:
\begin{equation}
f_{\rm vis} = \text{sinc}(\Delta\tau/\tau_*) \approx 1 - \frac{1}{6}\left(\frac{\Delta\tau}{\tau_*}\right)^2.
\end{equation}

\paragraph{Numerical value.}
With $\Delta\tau/\tau_* \sim 0.1$:
\begin{equation}
f_{\rm vis} \approx 1 - 0.02 = 0.98.
\end{equation}

\subsubsection{Doppler Factor $f_{\rm Dop}$}

The Doppler contribution from baryon velocity perturbations is:
\begin{equation}
\Theta_{\rm Dop} = \hat{n} \cdot \mathbf{v}_b,
\end{equation}
where $\hat{n}$ is the line-of-sight direction.

\paragraph{Effect on asymmetry.}
The Doppler term is $90^\circ$ out of phase with the acoustic term. When projected onto the line of sight and averaged, this reduces the effective asymmetry:
\begin{equation}
f_{\rm Dop} \approx 0.90.
\end{equation}

\subsubsection{Total Asymmetry}

Combining all factors:
\begin{equation}
A = 0.474 \times 0.50 \times 0.98 \times 0.90 = 0.209.
\end{equation}

%-----------------------------------------------------------------------------
\subsection{Peak Ratio Derivation}\label{app:peak-ratio}
%-----------------------------------------------------------------------------

\paragraph{Definition.}
The peak ratio is:
\begin{equation}
R \equiv \frac{H_1}{H_2} = \frac{(\text{first peak height})}{(\text{second peak height})}.
\end{equation}

\paragraph{Relation to asymmetry.}
For the angular power spectrum $C_\ell$, the peak heights scale as:
\begin{equation}
H_n \propto \left[(1 + (-1)^{n+1} A)\right]^2.
\end{equation}
Hence:
\begin{equation}
R = \frac{(1+A)^2}{(1-A)^2} = \left(\frac{1+A}{1-A}\right)^2.
\end{equation}

\paragraph{Result.}
With $A = 0.209$:
\begin{equation}
\boxed{R = \left(\frac{1.209}{0.791}\right)^2 = (1.528)^2 = 2.34}
\end{equation}

\paragraph{Comparison with observation.}
Planck measures $R \approx 2.4$. The agreement is within 2.5\%.

%-----------------------------------------------------------------------------
\subsection{Why the $1/\mu$ Enhancement Cancels}\label{app:mu-cancellation}
%-----------------------------------------------------------------------------

\paragraph{Key insight.}
In $\psi$-gravity, the driving term is enhanced: $\Phi_\psi = \Phi/\mu$. But this enhancement affects \textit{both} odd and even peaks equally.

\paragraph{Mathematical demonstration.}
The acoustic equation~\eqref{eq:app-acoustic} has driving term:
\begin{equation}
F(k) = -\frac{k^2}{1+R_b} \Phi_\psi = -\frac{k^2}{1+R_b} \frac{\Phi}{\mu}.
\end{equation}
The oscillation amplitude scales as:
\begin{equation}
|\Theta| \propto \frac{|F|}{c_s^2 k^2} \propto \frac{|\Phi|/\mu}{c_s^2} \propto \frac{1}{\mu}.
\end{equation}

\textit{All} peaks (odd and even) are enhanced by $1/\mu$. In the ratio:
\begin{equation}
R = \frac{H_1}{H_2} = \frac{|\Theta_{\rm odd}|^2}{|\Theta_{\rm even}|^2} \propto \frac{(1/\mu)^2}{(1/\mu)^2} = 1 \times (\text{baryon physics}).
\end{equation}

The $\mu$-enhancement drops out of the ratio. What survives is the baryon loading factor, which depends only on $R_b$---a quantity fixed by BBN and completely independent of dark matter.

\paragraph{Translation to $\Lambda$CDM language.}
In $\Lambda$CDM, the ``dark matter fraction'' $f_c = \Omega_c/(\Omega_c+\Omega_b) \approx 0.84$ enters the peak ratio. In DFD, this same number arises from:
\begin{equation}
f_{\rm DFD} = 1 - \mu_{\rm eff} \times (\text{projection factors}).
\end{equation}
There are no dark matter particles; $f_c$ is just another parameterization of $\mu(x)$ effects.

%-----------------------------------------------------------------------------
\subsection{$\psi$-Lensing and Peak Location}\label{app:psi-lensing}
%-----------------------------------------------------------------------------

\paragraph{The problem.}
Standard GR calculations without CDM give $\ell_1 \approx 297$, not the observed $\ell_1 \approx 220$. This has been cited as ``proof'' that dark matter is required.

\paragraph{The resolution.}
This argument assumes GR propagation with fixed $c$ and straight-line photon paths. In $\psi$-physics, light travels through a medium with varying refractive index $n = e^\psi$, producing gradient-index (GRIN) optics effects.

\subsubsection{Gradient-Index Optics}

\paragraph{Basic physics.}
In a medium with spatially varying $n(\mathbf{x})$, light rays follow curved paths according to Fermat's principle. For a gradient $\nabla n$, rays bend toward regions of higher $n$.

\paragraph{Angular magnification.}
For a GRIN lens with $n$ varying along the line of sight:
\begin{equation}
\frac{\theta_{\rm obs}}{\theta_{\rm emit}} = \frac{n_{\rm emit}}{n_{\rm obs}}.
\end{equation}

If $n_{\rm emit} > n_{\rm obs}$ (higher $n$ at source):
\begin{itemize}[nosep]
\item $\theta_{\rm obs} > \theta_{\rm emit}$: angular scales are \textit{magnified}
\item Observed $\ell$ is \textit{smaller} than ``true'' $\ell$ (since $\ell \propto 1/\theta$)
\end{itemize}

\subsubsection{Application to CMB}

\paragraph{$\psi$-gradient.}
With $n = e^\psi$, the angular scaling becomes:
\begin{equation}
\frac{\theta_{\rm obs}}{\theta_{\rm emit}} = e^{\psi_{\rm CMB} - \psi_{\rm here}} = e^{\Delta\psi}.
\end{equation}

\paragraph{Peak location relation.}
\begin{equation}
\ell_{\rm obs} = \ell_{\rm true} \times \frac{\theta_{\rm true}}{\theta_{\rm obs}} = \ell_{\rm true} \times e^{-\Delta\psi}.
\end{equation}

\paragraph{Required gradient.}
To obtain $\ell_{\rm obs} = 220$ from $\ell_{\rm true} = 297$:
\begin{align}
220 &= 297 \times e^{-\Delta\psi}, \\
e^{-\Delta\psi} &= 220/297 = 0.74, \\
\Delta\psi &= -\ln(0.74) = 0.30.
\end{align}

\paragraph{Physical interpretation.}
$\Delta\psi = \psi_{\rm CMB} - \psi_{\rm here} = 0.30$ means:
\begin{itemize}[nosep]
\item $\psi$ was $0.30$ higher at CMB than today
\item $n_{\rm CMB}/n_{\rm here} = e^{0.30} = 1.35$ (35\% higher refractive index)
\item $c_{\rm CMB}/c_{\rm here} = e^{-0.30} = 0.74$ (26\% slower light speed)
\end{itemize}

This is a \textit{modest} gradient---not fine-tuned.

%-----------------------------------------------------------------------------
\subsection{Consistency Checks}\label{app:cmb-consistency}
%-----------------------------------------------------------------------------

\paragraph{Self-consistency of $\Delta\psi = 0.30$.}

\begin{enumerate}
\item \textbf{$\alpha$-variation bounds.}
With $\alpha(\psi) = \alpha_0(1 + k_\alpha\psi)$ and $k_\alpha = \alpha^2/(2\pi) \approx 8.5 \times 10^{-6}$ (Sec.~\ref{sec:alpha-kalpha}):
\begin{equation}
\frac{\Delta\alpha}{\alpha} = k_\alpha \Delta\psi \approx 8.5 \times 10^{-6} \times 0.30 \approx 2.5 \times 10^{-6}.
\end{equation}
This is $\sim 2.5$ ppm---well within observational bounds. The quasar $\alpha$-variation literature constrains $|\Delta\alpha/\alpha| \lesssim 10^{-5}$ at $z \sim 2$--$3$, and CMB constraints are $|\Delta\alpha/\alpha| \lesssim 10^{-3}$. DFD satisfies both with ample margin.

\textit{Note:} The coupling $k_\alpha = \alpha^2/(2\pi)$ governs electromagnetic variation; this is distinct from the acceleration coupling $k_a = 3/(8\alpha) \approx 51$ that appears in galactic dynamics.

\item \textbf{BBN compatibility.}
BBN occurs at $T \sim 1$ MeV, much earlier than CMB ($T \sim 0.3$ eV). If $\psi$-evolution is monotonic, $\Delta\psi_{\rm BBN}$ could be larger, but BBN physics depends primarily on nuclear rates, not optical effects. The constraint is on $\alpha_{\rm BBN}$, which can accommodate $\mathcal{O}(10\%)$ variations.

\item \textbf{Late-time $\psi$.}
Today, $\psi_{\rm here} \equiv 0$ by convention. Local physics is unaffected by the absolute value of $\psi$---only gradients matter.
\end{enumerate}

%-----------------------------------------------------------------------------
\subsection{Comparison with $\Lambda$CDM}\label{app:lcdm-comparison}
%-----------------------------------------------------------------------------

\paragraph{Feature comparison between $\Lambda$CDM and $\psi$-Cosmology.}

\begin{center}
\renewcommand{\arraystretch}{1.3}%
\resizebox{\columnwidth}{!}{%
\begin{tabular}{lcc}
\toprule
Feature & $\Lambda$CDM & $\psi$-Cosmology \\
\midrule
Peak ratio $R$ & CDM-driven ($\Omega_c$) & Baryon loading ($R_b$) \\
Peak location $\ell_1$ & GR distances (with CDM) & $\psi$-lensing ($\Delta\psi$) \\
Free parameters & $\Omega_c$, $\Omega_\Lambda$, \ldots & None (locked from galaxies) \\
Dark matter & Particles (undetected) & $\mu(x)$ effect (no particles) \\
Dark energy & $\Lambda$ (unexplained) & Optical illusion \\
\bottomrule
\end{tabular}%
}
\end{center}

\paragraph{Key difference.}
$\Lambda$CDM introduces dark matter \textit{particles} to explain the CMB. DFD explains the same observations using $\psi$-physics:
\begin{itemize}[nosep]
\item Peak ratio: baryon loading (same $R_b$ from BBN)
\item Peak location: $\psi$-lensing (new effect from $n = e^\psi$)
\end{itemize}

There are no new particles, just new understanding of how light propagates in the $\psi$-universe.

%-----------------------------------------------------------------------------
\subsection{Falsifiable Predictions}\label{app:cmb-predictions}
%-----------------------------------------------------------------------------

The $\psi$-CMB solution makes specific predictions beyond the peak structure:

\begin{enumerate}
\item \textbf{Distance duality consistency.}
Etherington's reciprocity holds exactly in DFD's optical metric:
\begin{equation}
\frac{D_L}{(1+z)^2 D_A} = 1.
\end{equation}
Both $D_L$ and $D_A$ are screened equally by $e^{\Delta\psi_{\rm screen}}$, so the ratio cancels. Observational confirmation ($\eta = 1.01 \pm 0.02$) validates the metric structure. Any detected violation would falsify DFD's single-metric framework.

\item \textbf{Redshift-dependent $c_{\rm eff}$.}
If $c(\psi) = c_0 e^{-\psi}$ varies along the line of sight, time-of-arrival measurements for transient events at different redshifts could reveal this.

\item \textbf{Polarization consistency.}
The $\psi$-lensing should affect E-mode and B-mode polarization consistently. Any inconsistency would falsify the model.

\item \textbf{Higher peaks.}
The third peak ($\ell_3$) and beyond should follow the same $\psi$-lensing relation. If $\ell_3/\ell_1$ deviates from the predicted ratio, the model is ruled out.
\end{enumerate}

\paragraph{Ultimate test.}
If detailed numerical $\psi$-Boltzmann calculations show that peak ratio and peak location \textit{cannot} be simultaneously fit with a single consistent $\Delta\psi$, the $\psi$-CMB solution is falsified.
